\section{Stage 1: finalize the index and vertex streams}
\label{sec:finalize}

\subsection{One dispatch, two tails}

After the host reserves the spans, the finalize dispatch runs with
\begin{equation}
  \text{work} = \text{indexCount} + \text{nodeCount}.
\end{equation}
The first \code{indexCount} invocations write the index stream; the
remaining \code{nodeCount} invocations write the vertex payload:
\begin{lstlisting}[style=glsl,caption={Finalize stage.},label={lst:finalize}]
void finalizeStage() {
    uint gi = gl_GlobalInvocationID.x;
    if (gi < pc.indexCount) { writeIndex(gi); return; }
    uint node = gi - pc.indexCount;
    if (node >= pc.nodeCount) return;
    writeVertexForNode(node);
}
\end{lstlisting}
The index range is first because it is usually much larger than the node
count.  The two ranges are disjoint, so one dispatch covers both without a
barrier between them.

\subsection{Index remap}

For index $i$, the raw key is the node id:
\begin{equation}
  \text{slot} = \text{nodeSlot}(\text{rawKeys}[i]).
\end{equation}
If the slot is negative, the host's slot assignment and the walk disagree.
The shader increments \code{droppedCount} and fails the chunk rather than
writing slot 0.  Writing 0 would silently connect the corner to an
arbitrary vertex and produce a stretched triangle.

\subsection{Vertex reconstruction}

The vertex is rebuilt from the node, not from the walk's temporary cell:
\begin{enumerate}
\item QEF position from the node's eight corners and cube;
\item analytic normal at that position;
\item brush, falling back to the chunk's \code{fromBrush} when the node's
   brush is the discard sentinel;
\item HSV tint, texture type, and mapping data.
\end{enumerate}
This is why the raw key is enough.  The CPU walker keeps the cell payload
in the triangle, but that payload is a pure function of the node.  The GPU
recomputes it in stage 1, which also means the finalize stage does not need
to read the walk's scratch arrays.

\subsection{The index base and alignment}

The chunk's index span is bound with an offset.  The offset must be aligned
to the device's minimum storage-buffer offset alignment.  Since the element size is four
bytes and the alignment can be 256, the host aligns the descriptor offset
down and passes \code{relIndexBase} = the number of index elements between
the aligned offset and the first span element.  The shader writes
\begin{equation}
  \text{poolIdx}[\text{relIndexBase} + i] = \text{slot}.
\end{equation}
The descriptor range is
\begin{equation}
  (\text{relIndexBase} + \text{indexCapacity})\cdot4
\end{equation}
bytes, which includes the padding.  Without the extra padding, the last
elements of the span fall outside the bound range and are never written,
producing stretched triangles from stale pool data.

\subsection{The vertex base and alignment}

The vertex span is handled the same way with a 64-byte stride.  The offset
is aligned down; \code{relVertexBase} accounts for the 0--3 element
padding.  The shader writes
\begin{equation}
  \text{poolVerts}[(\text{relVertexBase} + s)\cdot4 + k].
\end{equation}
The range is
\begin{equation}
  (\text{relVertexBase} + \text{vertexCapacity})\cdot64.
\end{equation}
Because 64 is a multiple of 4 and the largest required alignment is 256, a
span start can be misaligned by at most three vertices.  The padding logic
is exact and costs at most 192 bytes per chunk.

\subsection{Drop accounting}

Every write that cannot fit increments \code{droppedCount}.  The host reads
the counter after the fence.  The chunk is considered valid only when
\begin{equation}
  \text{droppedCount}=0 \quad\land\quad
  \text{overflowCount}=0 \quad\land\quad
  \text{poolVertexCount}\le\text{vertexCapacity}.
\end{equation}
If any check fails, the reserved span is released and the chunk is skipped.
The next edit retries it.  The app never draws a partially written span.

\subsection{The raw-key buffer}

The raw-key buffer is device-local and sized to the key capacity.  It is
written in stage 0 and read in stage 1.  Its lifetime is one topology
submission (two dispatches separated by a host read).  The buffer is
reused across chunks; the host resets the counters before each walk.  The
key capacity is selected from a per-cell hint:
\begin{equation}
  \text{keys} = \min\bigl(\text{cellCount}\cdot25+64,\;2^{23}\bigr).
\end{equation}
A walk that exceeds the hint retries with twice the capacity, up to the
hard limit.  The measured worst chunk is about 3.8 million raw vertices,
below the 8 million limit.
